\documentclass[11pt,a4paper]{article}

%====================================================
% PACKAGES
%====================================================
\usepackage[utf8]{inputenc}
\usepackage[T1]{fontenc}
\usepackage[french]{babel}
\usepackage[a4paper,margin=1.8cm]{geometry}
\usepackage{amsmath,amssymb,mathtools}
\usepackage{lmodern}
\usepackage{xcolor}
\usepackage{enumitem}
\usepackage{fancyhdr}
\usepackage{lastpage}
\usepackage{array}
\usepackage{hyperref}
\usepackage{titlesec}

%====================================================
% MISE EN PAGE
%====================================================
\definecolor{bleuprepa}{RGB}{22,70,120}
\definecolor{grisclair}{RGB}{245,245,245}
\definecolor{grisfonce}{RGB}{80,80,80}

\pagestyle{fancy}
\fancyhf{}
\fancyhead[L]{\textbf{Algèbre linéaire}}
\fancyhead[C]{\textbf{TD -- Applications linéaires}}
\fancyhead[R]{\textbf{Page \thepage/\pageref{LastPage}}}
\setlength{\headheight}{15pt}

\hypersetup{
    colorlinks=true,
    linkcolor=bleuprepa,
    urlcolor=bleuprepa
}

\titleformat{\section}
  {\Large\bfseries\color{bleuprepa}}
  {\thesection}{0.7em}{}

\titleformat{\subsection}
  {\large\bfseries\color{grisfonce}}
  {\thesubsection}{0.7em}{}

\setlist[enumerate]{label=\textbf{\alph*)}, leftmargin=1.1cm, itemsep=0.08cm}
\setlist[itemize]{leftmargin=0.9cm, itemsep=0.08cm}

%====================================================
% COMMANDES
%====================================================
\newcounter{exo}
\newcommand{\exo}[3]{
    \refstepcounter{exo}
    \vspace{0.45cm}
    \noindent
    \colorbox{grisclair}{
    \parbox{0.965\linewidth}{
    \textbf{\textcolor{bleuprepa}{Exercice \theexo\ -- #1}}
    \hfill
    \textbf{[#2 \; #3]}
    }}
    \vspace{0.15cm}
}

\newcommand{\corr}[1]{
    \vspace{0.4cm}
    \noindent\textbf{\textcolor{bleuprepa}{Corrigé de l'exercice #1.}}
    \par\vspace{0.1cm}
}

\newcommand{\R}{\mathbb{R}}
\newcommand{\K}{\mathbb{K}}
\newcommand{\Ker}{\operatorname{Ker}}
\newcommand{\Img}{\operatorname{Im}}
\newcommand{\rg}{\operatorname{rg}}
\newcommand{\Mat}{\operatorname{Mat}}
\newcommand{\Vect}{\operatorname{Vect}}
\newcommand{\Id}{\operatorname{Id}}

%====================================================
% DOCUMENT
%====================================================
\begin{document}

\begin{center}
    {\Large \textbf{Classe préparatoire scientifique}}\\[0.2cm]
    {\Huge \textbf{\textcolor{bleuprepa}{TD complet -- Applications linéaires}}}\\[0.2cm]
    {\Large MPSI/PCSI/MP -- Algèbre linéaire}\\[0.25cm]
    \textbf{Nom :} \underline{\hspace{5cm}}
    \hfill
    \textbf{Date :} \underline{\hspace{3cm}}
\end{center}

\vspace{0.35cm}
\hrule
\vspace{0.25cm}

\noindent\textbf{Objectifs.}
Ce TD a pour but de maîtriser les applications linéaires, le calcul de noyau, d'image et de rang, ainsi que les critères d'injectivité et de surjectivité. Les exercices sont progressifs : application directe, méthode, démonstration, réflexion, puis oraux de concours.

\vspace{0.2cm}
\noindent\textbf{Types d'exercices.}
APP = application directe ; METH = méthode ; DEM = démonstration ; PREP = réflexion ; ORAL = type concours.

\vspace{0.25cm}
\hrule

\tableofcontents
\newpage

%====================================================
\section{Énoncés}
%====================================================

%====================================================
\subsection{Exercices d'application directe}
%====================================================

\exo{Vérifier la linéarité d'une application de $\R^2$ dans $\R^2$}{APP}{$\star$}
On considère
\[
u:\R^2\to\R^2,\qquad u(x,y)=(2x-y,\;x+3y).
\]
\begin{enumerate}
    \item Calculer $u(0,0)$ et rappeler pourquoi ce test est nécessaire mais non suffisant.
    \item Vérifier l'additivité de $u$ pour deux vecteurs quelconques $(x_1,y_1)$ et $(x_2,y_2)$.
    \item Vérifier l'homogénéité de $u$ pour un scalaire $\lambda\in\R$.
    \item Conclure rigoureusement que $u$ est linéaire.
\end{enumerate}

\exo{Reconnaître une application non linéaire}{APP}{$\star$}
On considère
\[
f:\R^2\to\R,\qquad f(x,y)=x^2+y.
\]
\begin{enumerate}
    \item Calculer $f(1,0)$, $f(2,0)$ et $f(0,0)$.
    \item Comparer $f((1,0)+(1,0))$ et $f(1,0)+f(1,0)$.
    \item Conclure que $f$ n'est pas linéaire.
    \item Expliquer en une phrase pourquoi la présence d'un terme quadratique est incompatible avec la linéarité.
\end{enumerate}

\exo{Forme linéaire et noyau dans $\R^3$}{APP}{$\star\star$}
On considère
\[
\varphi:\R^3\to\R,\qquad \varphi(x,y,z)=x-2y+3z.
\]
\begin{enumerate}
    \item Montrer que $\varphi$ est linéaire.
    \item Déterminer $\Ker(\varphi)$ sous forme paramétrique.
    \item Donner une base de $\Ker(\varphi)$.
    \item Donner la dimension du noyau et interpréter géométriquement le résultat.
\end{enumerate}

\exo{Noyau d'une application de $\R^2$ dans $\R$}{APP}{$\star\star$}
On considère
\[
u:\R^2\to\R,\qquad u(x,y)=x-y.
\]
\begin{enumerate}
    \item Résoudre l'équation $u(x,y)=0$.
    \item Décrire $\Ker(u)$ comme un sous-espace vectoriel.
    \item Donner une base de $\Ker(u)$.
    \item L'application $u$ est-elle injective ? Justifier avec le noyau.
\end{enumerate}

\exo{Image d'une application de $\R^3$ dans $\R^2$}{APP}{$\star\star$}
On considère
\[
u:\R^3\to\R^2,\qquad u(x,y,z)=(x+y,\;y+z).
\]
\begin{enumerate}
    \item Calculer $u(e_1)$, $u(e_2)$ et $u(e_3)$, où $(e_1,e_2,e_3)$ est la base canonique.
    \item En déduire une famille génératrice de $\Img(u)$.
    \item Extraire une base de $\Img(u)$.
    \item Calculer $\rg(u)$.
\end{enumerate}

\exo{Injectivité et surjectivité}{APP}{$\star\star$}
On reprend
\[
u:\R^3\to\R^2,\qquad u(x,y,z)=(x+y,\;y+z).
\]
\begin{enumerate}
    \item Déterminer $\Ker(u)$.
    \item L'application est-elle injective ?
    \item L'application est-elle surjective ?
    \item Vérifier le théorème du rang.
\end{enumerate}

\exo{Matrice d'une application dans les bases canoniques}{APP}{$\star\star$}
Soit
\[
u:\R^2\to\R^3,\qquad u(x,y)=(x+y,\;2x-y,\;3y).
\]
\begin{enumerate}
    \item Calculer $u(1,0)$ et $u(0,1)$.
    \item Construire la matrice de $u$ dans les bases canoniques.
    \item Vérifier que le produit matrice-colonne redonne bien $u(x,y)$.
    \item Déterminer si $u$ peut être surjective.
\end{enumerate}

\exo{Dérivation sur $\R_2[X]$}{APP}{$\star\star$}
On considère
\[
D:\R_2[X]\to\R_2[X],\qquad P\mapsto P'.
\]
On travaille dans la base $B=(1,X,X^2)$.
\begin{enumerate}
    \item Calculer $D(1)$, $D(X)$ et $D(X^2)$.
    \item Écrire la matrice de $D$ dans la base $B$.
    \item Déterminer $\Ker(D)$.
    \item Déterminer $\Img(D)$.
    \item Vérifier le théorème du rang.
\end{enumerate}

\exo{Application donnée par une matrice}{APP}{$\star\star$}
Soit
\[
A=\begin{pmatrix}1&2\\2&4\end{pmatrix}
\]
et soit $u$ l'application linéaire de $\R^2$ dans $\R^2$ de matrice $A$ dans la base canonique.
\begin{enumerate}
    \item Déterminer explicitement $u(x,y)$.
    \item Calculer $\Ker(u)$.
    \item Calculer $\Img(u)$.
    \item Donner $\rg(u)$.
    \item L'application est-elle un automorphisme de $\R^2$ ?
\end{enumerate}

\exo{Projecteur canonique de $\R^3$}{APP}{$\star\star$}
On considère
\[
p:\R^3\to\R^3,\qquad p(x,y,z)=(x,y,0).
\]
\begin{enumerate}
    \item Montrer que $p$ est linéaire.
    \item Calculer $p^2$.
    \item Déterminer $\Ker(p)$.
    \item Déterminer $\Img(p)$.
    \item Interpréter géométriquement $p$.
\end{enumerate}

%====================================================
\subsection{Exercices de méthode}
%====================================================

\exo{Construire une application par l'image d'une base}{METH}{$\star\star$}
On considère la base canonique $(e_1,e_2,e_3)$ de $\R^3$. On cherche une application linéaire
\[
u:\R^3\to\R^2
\]
telle que
\[
u(e_1)=(1,0),\qquad u(e_2)=(0,1),\qquad u(e_3)=(1,1).
\]
\begin{enumerate}
    \item Justifier qu'une telle application existe et est unique.
    \item Déterminer l'expression de $u(x,y,z)$.
    \item Écrire la matrice de $u$ dans les bases canoniques.
    \item Calculer $\Ker(u)$.
    \item Calculer $\Img(u)$ et $\rg(u)$.
\end{enumerate}

\exo{Application paramétrée}{METH}{$\star\star\star$}
Pour $\lambda\in\R$, on définit
\[
u_\lambda:\R^2\to\R^2,\qquad u_\lambda(x,y)=(x+\lambda y,\;\lambda x+y).
\]
\begin{enumerate}
    \item Montrer que $u_\lambda$ est linéaire.
    \item Écrire sa matrice $A_\lambda$.
    \item Résoudre le système $u_\lambda(x,y)=(0,0)$ selon les valeurs de $\lambda$.
    \item Pour quelles valeurs de $\lambda$ l'application est-elle injective ?
    \item Pour quelles valeurs de $\lambda$ l'application est-elle surjective ?
    \item Étudier précisément les cas $\lambda=1$ et $\lambda=-1$.
\end{enumerate}

\exo{Rang, noyau et image}{METH}{$\star\star\star$}
On considère
\[
u:\R^3\to\R^2,\qquad u(x,y,z)=(x+y+z,\;x+2y+3z).
\]
\begin{enumerate}
    \item Vérifier que $u$ est linéaire.
    \item Calculer les images des vecteurs de la base canonique.
    \item Déterminer une base de $\Img(u)$.
    \item Déterminer $\Ker(u)$.
    \item Vérifier le théorème du rang.
\end{enumerate}

\exo{Image d'une base et rang}{METH}{$\star\star\star$}
Soit $u:E\to F$ linéaire, où $E$ est de dimension $3$, et soit $(e_1,e_2,e_3)$ une base de $E$.
On suppose que
\[
u(e_1)=f_1,\qquad u(e_2)=f_2,\qquad u(e_3)=f_1+f_2.
\]
\begin{enumerate}
    \item Montrer que $\Img(u)=\Vect(f_1,f_2)$.
    \item À quelle condition sur $(f_1,f_2)$ a-t-on $\rg(u)=2$ ?
    \item Dans ce cas, calculer $\dim(\Ker(u))$.
    \item Donner un vecteur non nul de $\Ker(u)$ en fonction de $e_1,e_2,e_3$.
\end{enumerate}

\exo{Étude matricielle complète sans déterminant}{METH}{$\star\star\star$}
Soit
\[
A=\begin{pmatrix}
1&1&0\\
0&1&1\\
1&0&1
\end{pmatrix}.
\]
On note $u$ l'endomorphisme de $\R^3$ de matrice $A$ dans la base canonique.
\begin{enumerate}
    \item Écrire $u(x,y,z)$.
    \item Résoudre $u(x,y,z)=(0,0,0)$.
    \item En déduire $\Ker(u)$.
    \item En déduire l'injectivité puis le rang.
    \item Déterminer $\Img(u)$.
\end{enumerate}

\exo{Application à valeurs dans un espace de polynômes}{METH}{$\star\star\star$}
On définit
\[
u:\R^2\to\R_2[X],\qquad u(a,b)=a(1+X)+b(X+X^2).
\]
\begin{enumerate}
    \item Montrer que $u$ est linéaire.
    \item Écrire la matrice de $u$ dans la base canonique de $\R^2$ et la base $(1,X,X^2)$.
    \item Déterminer $\Ker(u)$.
    \item Déterminer $\Img(u)$.
    \item L'application est-elle injective ? surjective ?
\end{enumerate}

\exo{Forme linéaire et hyperplan}{METH}{$\star\star\star$}
On considère
\[
\varphi:\R^4\to\R,\qquad
\varphi(x,y,z,t)=x+2y-z+t.
\]
\begin{enumerate}
    \item Montrer que $\varphi$ est une forme linéaire non nulle.
    \item Déterminer $\Ker(\varphi)$ sous forme paramétrique.
    \item Donner une base de $\Ker(\varphi)$.
    \item Vérifier que $\Ker(\varphi)$ est un hyperplan de $\R^4$.
\end{enumerate}

\exo{Image d'une famille génératrice}{METH}{$\star\star\star$}
Soit $u:E\to F$ linéaire et soit $(e_1,\dots,e_n)$ une famille génératrice de $E$.
\begin{enumerate}
    \item Soit $y\in\Img(u)$. Écrire $y=u(x)$ pour un certain $x\in E$.
    \item Décomposer $x$ comme combinaison linéaire des $e_i$.
    \item En déduire que $(u(e_1),\dots,u(e_n))$ engendre $\Img(u)$.
    \item Donner un exemple où cette famille n'est pas libre.
\end{enumerate}

\exo{Image d'une famille libre par une application injective}{METH}{$\star\star\star$}
Soit $u:E\to F$ linéaire injective.
\begin{enumerate}
    \item Soit $(e_1,\dots,e_n)$ une famille libre de $E$.
    \item Supposer qu'une relation
    \[
    \lambda_1u(e_1)+\cdots+\lambda_nu(e_n)=0
    \]
    est vérifiée.
    \item Utiliser la linéarité et l'injectivité de $u$.
    \item Conclure que $(u(e_1),\dots,u(e_n))$ est libre.
\end{enumerate}

\exo{Composition et matrices}{METH}{$\star\star\star$}
On considère
\[
u:\R^2\to\R^3,\qquad u(x,y)=(x+y,\;x-y,\;2y),
\]
et
\[
v:\R^3\to\R^2,\qquad v(a,b,c)=(a+c,\;b-c).
\]
\begin{enumerate}
    \item Écrire les matrices $A$ et $B$ de $u$ et $v$ dans les bases canoniques.
    \item Calculer $v\circ u$ directement.
    \item Calculer la matrice de $v\circ u$ par produit matriciel.
    \item Vérifier la cohérence des deux méthodes.
\end{enumerate}

%====================================================
\subsection{Exercices de démonstration}
%====================================================

\exo{Image du vecteur nul}{DEM}{$\star\star$}
Soit $u:E\to F$ linéaire.
\begin{enumerate}
    \item Écrire $0_E=0_E+0_E$.
    \item Appliquer $u$ à cette égalité.
    \item En déduire que $u(0_E)=0_F$.
\end{enumerate}

\exo{Le noyau est un sous-espace vectoriel}{DEM}{$\star\star$}
Soit $u:E\to F$ linéaire.
\begin{enumerate}
    \item Montrer que $0_E\in\Ker(u)$.
    \item Montrer la stabilité de $\Ker(u)$ par addition.
    \item Montrer la stabilité de $\Ker(u)$ par multiplication scalaire.
    \item Conclure.
\end{enumerate}

\exo{L'image est un sous-espace vectoriel}{DEM}{$\star\star$}
Soit $u:E\to F$ linéaire.
\begin{enumerate}
    \item Montrer que $0_F\in\Img(u)$.
    \item Soient $y_1,y_2\in\Img(u)$. Écrire $y_1=u(x_1)$ et $y_2=u(x_2)$.
    \item Montrer que $\lambda y_1+\mu y_2\in\Img(u)$.
    \item Conclure.
\end{enumerate}

\exo{Critère d'injectivité par le noyau}{DEM}{$\star\star\star$}
Soit $u:E\to F$ linéaire.
\begin{enumerate}
    \item Montrer que si $u$ est injective, alors $\Ker(u)=\{0_E\}$.
    \item Réciproquement, supposer $\Ker(u)=\{0_E\}$.
    \item Montrer que si $u(x)=u(y)$, alors $x-y\in\Ker(u)$.
    \item Conclure.
\end{enumerate}

\exo{Détermination par l'image d'une base}{DEM}{$\star\star\star$}
Soit $(e_1,\dots,e_n)$ une base de $E$ et soient $f_1,\dots,f_n$ des vecteurs de $F$.
\begin{enumerate}
    \item Montrer que tout $x\in E$ s'écrit de manière unique sous la forme
    \[
    x=\lambda_1e_1+\cdots+\lambda_ne_n.
    \]
    \item Définir $u(x)=\lambda_1f_1+\cdots+\lambda_nf_n$.
    \item Montrer que $u$ est linéaire.
    \item Montrer l'unicité d'une telle application.
\end{enumerate}

\exo{Théorème du rang}{DEM}{$\star\star\star\star$}
Soit $u:E\to F$ linéaire, avec $E$ de dimension finie.
\begin{enumerate}
    \item Choisir une base $(e_1,\dots,e_k)$ de $\Ker(u)$.
    \item Compléter cette famille en une base $(e_1,\dots,e_k,e_{k+1},\dots,e_n)$ de $E$.
    \item Montrer que $(u(e_{k+1}),\dots,u(e_n))$ engendre $\Img(u)$.
    \item Montrer que cette famille est libre.
    \item En déduire
    \[
    \dim(E)=\dim(\Ker(u))+\rg(u).
    \]
\end{enumerate}

\exo{Dimension de $\mathcal{L}(E,F)$}{DEM}{$\star\star\star\star$}
Soient $E$ et $F$ deux espaces vectoriels de dimensions respectives $n$ et $p$.
\begin{enumerate}
    \item Expliquer pourquoi une application linéaire $u:E\to F$ est déterminée par les images d'une base de $E$.
    \item Pour chaque vecteur de base de $E$, combien de coefficients faut-il pour décrire son image dans une base de $F$ ?
    \item En déduire que
    \[
    \dim(\mathcal{L}(E,F))=np.
    \]
\end{enumerate}

\exo{Composition d'applications linéaires}{DEM}{$\star\star\star$}
Soient $u:E\to F$ et $v:F\to G$ deux applications linéaires.
\begin{enumerate}
    \item Montrer que $v\circ u$ est bien définie de $E$ dans $G$.
    \item Vérifier que $(v\circ u)(\lambda x+\mu y)=\lambda(v\circ u)(x)+\mu(v\circ u)(y)$.
    \item Conclure.
\end{enumerate}

\exo{Inverse d'un isomorphisme}{DEM}{$\star\star\star\star$}
Soit $u:E\to F$ un isomorphisme.
\begin{enumerate}
    \item Soient $y_1,y_2\in F$. Poser $x_1=u^{-1}(y_1)$ et $x_2=u^{-1}(y_2)$.
    \item Calculer $u(\lambda x_1+\mu x_2)$.
    \item En déduire une expression de $u^{-1}(\lambda y_1+\mu y_2)$.
    \item Conclure que $u^{-1}$ est linéaire.
\end{enumerate}

\exo{Forme linéaire non nulle et hyperplan}{DEM}{$\star\star\star\star$}
Soit $E$ de dimension finie $n$ et soit $\varphi:E\to\R$ une forme linéaire non nulle.
\begin{enumerate}
    \item Montrer que $\Img(\varphi)=\R$.
    \item Appliquer le théorème du rang.
    \item En déduire que $\dim(\Ker(\varphi))=n-1$.
    \item Conclure que $\Ker(\varphi)$ est un hyperplan.
\end{enumerate}

%====================================================
\subsection{Exercices de réflexion}
%====================================================

\exo{Applications injectives non surjectives}{PREP}{$\star\star\star$}
\begin{enumerate}
    \item Peut-on construire une application linéaire injective de $\R^2$ dans $\R^3$ ?
    \item En donner un exemple.
    \item Peut-elle être surjective ?
    \item Justifier par un argument de dimension.
\end{enumerate}

\exo{Applications surjectives non injectives}{PREP}{$\star\star\star$}
\begin{enumerate}
    \item Peut-on construire une application linéaire surjective de $\R^3$ dans $\R^2$ ?
    \item En donner un exemple.
    \item Peut-elle être injective ?
    \item Justifier par le théorème du rang.
\end{enumerate}

\exo{Endomorphismes de $\R$}{PREP}{$\star\star$}
On considère une application linéaire $u:\R\to\R$.
\begin{enumerate}
    \item Poser $a=u(1)$.
    \item Montrer que pour tout $x\in\R$, $u(x)=ax$.
    \item En déduire toutes les applications linéaires de $\R$ dans $\R$.
\end{enumerate}

\exo{Endomorphismes de rang $1$}{PREP}{$\star\star\star\star$}
Soit $u:E\to E$ un endomorphisme de rang $1$.
\begin{enumerate}
    \item Montrer que $\Img(u)$ est une droite vectorielle.
    \item Soit $e\neq 0$ un vecteur de $\Img(u)$. Montrer que pour tout $x\in E$, il existe un scalaire $\varphi(x)$ tel que
    \[
    u(x)=\varphi(x)e.
    \]
    \item Montrer que $\varphi$ est une forme linéaire.
    \item Décrire $\Ker(u)$ en fonction de $\varphi$.
\end{enumerate}

\exo{Projecteurs abstraits}{PREP}{$\star\star\star\star$}
Soit $p:E\to E$ tel que $p^2=p$.
\begin{enumerate}
    \item Montrer que pour tout $x\in E$,
    \[
    x=(x-p(x))+p(x).
    \]
    \item Montrer que $x-p(x)\in\Ker(p)$.
    \item Montrer que $p(x)\in\Img(p)$.
    \item Montrer que
    \[
    E=\Ker(p)\oplus\Img(p).
    \]
\end{enumerate}

\exo{Symétries abstraites}{PREP}{$\star\star\star\star$}
Soit $s:E\to E$ tel que $s^2=\Id_E$.
\begin{enumerate}
    \item Montrer que $(s-\Id_E)(s+\Id_E)=0$.
    \item Pour $x\in E$, montrer que
    \[
    x=\frac{x+s(x)}{2}+\frac{x-s(x)}{2}.
    \]
    \item Montrer que la première composante appartient à $\Ker(s-\Id_E)$.
    \item Montrer que la seconde appartient à $\Ker(s+\Id_E)$.
    \item Conclure que
    \[
    E=\Ker(s-\Id_E)\oplus\Ker(s+\Id_E).
    \]
\end{enumerate}

\exo{Nilpotence d'ordre $2$}{PREP}{$\star\star\star\star$}
Soit $u:E\to E$ tel que $u^2=0$.
\begin{enumerate}
    \item Montrer que $\Img(u)\subset\Ker(u)$.
    \item En déduire que $\rg(u)\leq \dim(\Ker(u))$.
    \item En utilisant le théorème du rang, montrer que
    \[
    2\rg(u)\leq \dim(E).
    \]
\end{enumerate}

\exo{Applications en dimension égale}{PREP}{$\star\star\star$}
Soit $u:E\to F$ linéaire, avec $\dim(E)=\dim(F)<+\infty$.
\begin{enumerate}
    \item Montrer que si $u$ est injective, alors $u$ est surjective.
    \item Montrer que si $u$ est surjective, alors $u$ est injective.
    \item En déduire trois caractérisations équivalentes d'un isomorphisme.
\end{enumerate}

\exo{Base adaptée à un projecteur}{PREP}{$\star\star\star\star$}
Soit $p:E\to E$ un projecteur.
\begin{enumerate}
    \item Choisir une base de $\Ker(p)$ et une base de $\Img(p)$.
    \item Justifier que la réunion de ces deux bases est une base de $E$.
    \item Écrire la matrice de $p$ dans cette base adaptée.
\end{enumerate}

\exo{Rang d'une composée}{PREP}{$\star\star\star\star$}
Soient $u:E\to F$ et $v:F\to G$ linéaires, avec espaces de dimensions finies.
\begin{enumerate}
    \item Montrer que $\Img(v\circ u)\subset \Img(v)$.
    \item Montrer que $\rg(v\circ u)\leq \rg(v)$.
    \item Montrer que $\Img(v\circ u)=v(\Img(u))$.
    \item En déduire que $\rg(v\circ u)\leq \rg(u)$.
\end{enumerate}

%====================================================
\subsection{Exercices type oraux Mines, Centrale, CCP-INP}
%====================================================

\exo{Oral CCP -- Étude complète d'une application}{ORAL}{$\star\star\star$}
On considère
\[
u:\R^3\to\R^3,\qquad u(x,y,z)=(x+y,\;y+z,\;x+z).
\]
\begin{enumerate}
    \item Vérifier que $u$ est linéaire.
    \item Déterminer $\Ker(u)$.
    \item En déduire si $u$ est injective.
    \item Déterminer $\Img(u)$.
    \item En déduire si $u$ est surjective.
    \item Déterminer si $\Img(u)$ est une droite, un plan ou tout $\R^3$.
\end{enumerate}

\exo{Oral Mines -- Projecteur}{ORAL}{$\star\star\star\star$}
Soit $p:E\to E$ tel que $p^2=p$.
\begin{enumerate}
    \item Montrer que $p$ fixe tous les vecteurs de $\Img(p)$.
    \item Montrer que $E=\Ker(p)\oplus\Img(p)$.
    \item Réciproquement, si $E=F\oplus G$, définir le projecteur sur $F$ parallèlement à $G$.
    \item Donner sa matrice dans une base adaptée.
\end{enumerate}

\exo{Oral Centrale -- Symétrie}{ORAL}{$\star\star\star\star$}
Soit $s:E\to E$ tel que $s^2=\Id_E$.
\begin{enumerate}
    \item Montrer que $s$ est bijective.
    \item Décomposer tout vecteur $x$ en somme d'un vecteur fixé par $s$ et d'un vecteur envoyé sur son opposé.
    \item En déduire une décomposition directe de $E$.
    \item Interpréter géométriquement.
\end{enumerate}

\exo{Oral Mines -- Rang $1$}{ORAL}{$\star\star\star\star\star$}
Soit $u:E\to F$ une application linéaire de rang $1$.
\begin{enumerate}
    \item Montrer que $\Img(u)$ est une droite.
    \item Montrer qu'il existe une forme linéaire $\varphi$ et un vecteur non nul $f\in F$ tels que
    \[
    u(x)=\varphi(x)f.
    \]
    \item Étudier le noyau de $u$.
    \item Réciproquement, donner une condition pour qu'une application de la forme $x\mapsto\varphi(x)f$ soit de rang $1$.
\end{enumerate}

\exo{Oral CCP -- Noyau et image d'une matrice}{ORAL}{$\star\star\star$}
Soit
\[
A=\begin{pmatrix}
1&2&1\\
2&4&2\\
1&1&0
\end{pmatrix}.
\]
\begin{enumerate}
    \item Calculer le noyau de l'application associée.
    \item Déterminer son image.
    \item Calculer son rang.
    \item Vérifier le théorème du rang.
\end{enumerate}

\exo{Oral Centrale -- Endomorphisme vérifiant $u^2=0$}{ORAL}{$\star\star\star\star$}
Soit $u:E\to E$ tel que $u^2=0$.
\begin{enumerate}
    \item Montrer que $\Img(u)\subset\Ker(u)$.
    \item Donner une conséquence sur le rang de $u$.
    \item Que peut-on dire si $\dim(E)=3$ ?
    \item Donner un exemple non nul dans $\R^2$.
\end{enumerate}

\exo{Oral Mines -- Matrice dans une base adaptée}{ORAL}{$\star\star\star\star$}
Dans $\R^2$, on définit
\[
u(x,y)=(3x+y,\;x+3y).
\]
On pose
\[
v_1=(1,1),\qquad v_2=(1,-1).
\]
\begin{enumerate}
    \item Montrer que $(v_1,v_2)$ est une base de $\R^2$.
    \item Calculer $u(v_1)$ et $u(v_2)$.
    \item En déduire la matrice de $u$ dans la base $(v_1,v_2)$.
    \item Comparer avec la matrice de $u$ dans la base canonique.
\end{enumerate}

\exo{Oral CCP -- Forme linéaire et hyperplan}{ORAL}{$\star\star\star$}
On considère
\[
H=\{(x,y,z,t)\in\R^4\mid x-y+2z+t=0\}.
\]
\begin{enumerate}
    \item Montrer que $H$ est le noyau d'une forme linéaire.
    \item Donner une base de $H$.
    \item Montrer que $H$ est un hyperplan.
\end{enumerate}

\exo{Oral Centrale -- Composition et noyau}{ORAL}{$\star\star\star\star$}
Soient $u:E\to F$ et $v:F\to G$ deux applications linéaires.
\begin{enumerate}
    \item Montrer que $\Ker(u)\subset\Ker(v\circ u)$.
    \item Donner une condition suffisante pour avoir égalité.
    \item Montrer que $\Img(v\circ u)\subset\Img(v)$.
    \item Donner un exemple où l'inclusion est stricte.
\end{enumerate}

\exo{Oral Mines -- Isomorphisme induit}{ORAL}{$\star\star\star\star\star$}
Soit $u:E\to F$ linéaire. On suppose que $E$ est de dimension finie.
\begin{enumerate}
    \item Choisir un supplémentaire $G$ de $\Ker(u)$ dans $E$.
    \item Montrer que la restriction de $u$ à $G$ est injective.
    \item Montrer que $u(G)=\Img(u)$.
    \item En déduire que $G$ est isomorphe à $\Img(u)$.
    \item Retrouver le théorème du rang.
\end{enumerate}

\newpage

%====================================================
\section{Corrigé complet}
%====================================================

\subsection{Corrigé des exercices d'application directe}

\corr{1}
On calcule d'abord
\[
u(0,0)=(0,0).
\]
Soient $X=(x_1,y_1)$ et $Y=(x_2,y_2)$. Alors
\[
u(X+Y)=u(x_1+x_2,y_1+y_2)
=(2x_1+2x_2-y_1-y_2,\;x_1+x_2+3y_1+3y_2).
\]
D'autre part,
\[
u(X)+u(Y)=(2x_1-y_1,x_1+3y_1)+(2x_2-y_2,x_2+3y_2),
\]
ce qui donne le même résultat. Enfin,
\[
u(\lambda x,\lambda y)=(2\lambda x-\lambda y,\lambda x+3\lambda y)=\lambda u(x,y).
\]
Donc $u$ est linéaire.

\corr{2}
On a
\[
f(1,0)=1,\qquad f(2,0)=4.
\]
Mais
\[
f((1,0)+(1,0))=f(2,0)=4,
\]
alors que
\[
f(1,0)+f(1,0)=1+1=2.
\]
L'additivité est fausse, donc $f$ n'est pas linéaire. La présence du terme $x^2$ explique ce défaut : l'application ne respecte pas les combinaisons linéaires.

\corr{3}
La linéarité vient du fait que $\varphi$ est combinaison linéaire des coordonnées :
\[
\varphi(\lambda X+\mu Y)=\lambda\varphi(X)+\mu\varphi(Y).
\]
Le noyau est donné par
\[
x-2y+3z=0.
\]
On pose $y=s$ et $z=t$, donc $x=2s-3t$. Ainsi
\[
(x,y,z)=s(2,1,0)+t(-3,0,1).
\]
Donc
\[
\Ker(\varphi)=\Vect((2,1,0),(-3,0,1)).
\]
Ces deux vecteurs sont non colinéaires, donc forment une base du noyau. La dimension vaut $2$. Géométriquement, le noyau est un plan vectoriel de $\R^3$.

\corr{4}
On résout
\[
x-y=0,
\]
donc $x=y$. Ainsi
\[
\Ker(u)=\{(t,t)\mid t\in\R\}=\Vect((1,1)).
\]
Le noyau est une droite vectorielle. Il n'est donc pas réduit à $\{0\}$. Par conséquent, l'application n'est pas injective.

\corr{5}
On calcule
\[
u(e_1)=u(1,0,0)=(1,0),\qquad
u(e_2)=u(0,1,0)=(1,1),\qquad
u(e_3)=u(0,0,1)=(0,1).
\]
Donc
\[
\Img(u)=\Vect((1,0),(1,1),(0,1)).
\]
Comme $(1,0)$ et $(0,1)$ forment une base de $\R^2$, on obtient
\[
\Img(u)=\R^2.
\]
Donc
\[
\rg(u)=2.
\]

\corr{6}
On cherche les solutions de
\[
x+y=0,\qquad y+z=0.
\]
On obtient $x=-y$ et $z=-y$, donc
\[
(x,y,z)=y(-1,1,-1).
\]
Ainsi
\[
\Ker(u)=\Vect((-1,1,-1)).
\]
Le noyau n'est pas nul, donc $u$ n'est pas injective. D'après l'exercice précédent, l'image est $\R^2$, donc $u$ est surjective. Enfin :
\[
3=\dim(\Ker(u))+\rg(u)=1+2.
\]

\corr{7}
On calcule
\[
u(1,0)=(1,2,0),\qquad u(0,1)=(1,-1,3).
\]
La matrice de $u$ est donc
\[
\Mat(u)=
\begin{pmatrix}
1&1\\
2&-1\\
0&3
\end{pmatrix}.
\]
Alors
\[
\begin{pmatrix}
1&1\\
2&-1\\
0&3
\end{pmatrix}
\begin{pmatrix}x\\y\end{pmatrix}
=
\begin{pmatrix}
x+y\\
2x-y\\
3y
\end{pmatrix}.
\]
L'application ne peut pas être surjective de $\R^2$ vers $\R^3$, car son rang est au plus $2$.

\corr{8}
On a
\[
D(1)=0,\qquad D(X)=1,\qquad D(X^2)=2X.
\]
Dans la base $B=(1,X,X^2)$ :
\[
[0]_B=\begin{pmatrix}0\\0\\0\end{pmatrix},\qquad
[1]_B=\begin{pmatrix}1\\0\\0\end{pmatrix},\qquad
[2X]_B=\begin{pmatrix}0\\2\\0\end{pmatrix}.
\]
Donc
\[
\Mat_B(D)=
\begin{pmatrix}
0&1&0\\
0&0&2\\
0&0&0
\end{pmatrix}.
\]
Le noyau est l'ensemble des polynômes constants :
\[
\Ker(D)=\Vect(1).
\]
L'image est
\[
\Img(D)=\Vect(1,X).
\]
Ainsi
\[
3=\dim(\Ker(D))+\rg(D)=1+2.
\]

\corr{9}
La matrice donne
\[
u(x,y)=(x+2y,2x+4y).
\]
Le noyau est donné par
\[
x+2y=0,
\]
la seconde équation étant la même multipliée par $2$. Donc
\[
\Ker(u)=\Vect((-2,1)).
\]
L'image est engendrée par les colonnes :
\[
\Img(u)=\Vect((1,2),(2,4))=\Vect((1,2)).
\]
Donc
\[
\rg(u)=1.
\]
L'application n'est pas un automorphisme, car son noyau n'est pas nul.

\corr{10}
L'application est linéaire car chaque coordonnée est combinaison linéaire des coordonnées de départ. On a
\[
p^2(x,y,z)=p(x,y,0)=(x,y,0)=p(x,y,z).
\]
Donc $p^2=p$, c'est un projecteur. On a
\[
\Ker(p)=\{(0,0,z)\mid z\in\R\}=\Vect((0,0,1)),
\]
et
\[
\Img(p)=\{(x,y,0)\mid x,y\in\R\}=\Vect((1,0,0),(0,1,0)).
\]
Géométriquement, $p$ est la projection sur le plan $(Oxy)$ parallèlement à l'axe $(Oz)$.

\subsection{Corrigé des exercices de méthode}

\corr{11}
Une application linéaire est entièrement déterminée par l'image d'une base. Il existe donc une unique application linéaire vérifiant ces conditions. Pour $(x,y,z)=xe_1+ye_2+ze_3$ :
\[
u(x,y,z)=xu(e_1)+yu(e_2)+zu(e_3).
\]
Donc
\[
u(x,y,z)=x(1,0)+y(0,1)+z(1,1)=(x+z,y+z).
\]
La matrice est
\[
A=
\begin{pmatrix}
1&0&1\\
0&1&1
\end{pmatrix}.
\]
Le noyau est défini par
\[
x+z=0,\qquad y+z=0.
\]
Donc $z=-x$ et $y=x$, d'où
\[
\Ker(u)=\Vect((1,1,-1)).
\]
Les colonnes $(1,0)$ et $(0,1)$ sont dans l'image, donc
\[
\Img(u)=\R^2,\qquad \rg(u)=2.
\]

\corr{12}
La matrice est
\[
A_\lambda=
\begin{pmatrix}
1&\lambda\\
\lambda&1
\end{pmatrix}.
\]
On résout
\[
\begin{cases}
x+\lambda y=0,\\
\lambda x+y=0.
\end{cases}
\]
Si $\lambda\neq 1$ et $\lambda\neq -1$, la seule solution est $(0,0)$, donc le noyau est nul. L'application est alors injective. Comme elle va de $\R^2$ dans $\R^2$, elle est aussi surjective.

Si $\lambda=1$, on obtient
\[
x+y=0,
\]
donc
\[
\Ker(u_1)=\Vect((1,-1)).
\]
Si $\lambda=-1$, on obtient
\[
x-y=0,
\]
donc
\[
\Ker(u_{-1})=\Vect((1,1)).
\]
Dans ces deux cas, l'application n'est pas injective et son rang vaut $1$.

\corr{13}
On calcule
\[
u(e_1)=(1,1),\qquad u(e_2)=(1,2),\qquad u(e_3)=(1,3).
\]
Les vecteurs $(1,1)$ et $(1,2)$ ne sont pas colinéaires, donc ils forment une base de $\R^2$. Ainsi
\[
\Img(u)=\R^2,\qquad \rg(u)=2.
\]
Le noyau est donné par
\[
\begin{cases}
x+y+z=0,\\
x+2y+3z=0.
\end{cases}
\]
En soustrayant la première équation à la seconde :
\[
y+2z=0,
\]
donc $y=-2z$. Puis
\[
x-2z+z=0,
\]
donc $x=z$. Ainsi
\[
(x,y,z)=z(1,-2,1).
\]
Donc
\[
\Ker(u)=\Vect((1,-2,1)).
\]
On vérifie bien
\[
3=1+2.
\]

\corr{14}
Comme $(e_1,e_2,e_3)$ est une base de $E$, l'image de $u$ est engendrée par
\[
u(e_1),u(e_2),u(e_3).
\]
Or
\[
u(e_3)=f_1+f_2.
\]
Donc
\[
\Img(u)=\Vect(f_1,f_2,f_1+f_2)=\Vect(f_1,f_2).
\]
Si $(f_1,f_2)$ est libre, alors $\rg(u)=2$. Dans ce cas, comme $\dim(E)=3$, le théorème du rang donne
\[
\dim(\Ker(u))=3-2=1.
\]
Enfin
\[
u(e_3-e_1-e_2)=u(e_3)-u(e_1)-u(e_2)=f_1+f_2-f_1-f_2=0.
\]
Donc
\[
e_3-e_1-e_2\in\Ker(u).
\]

\corr{15}
On a
\[
u(x,y,z)=(x+y,\;y+z,\;x+z).
\]
Le noyau est donné par
\[
\begin{cases}
x+y=0,\\
y+z=0,\\
x+z=0.
\end{cases}
\]
De $x+y=0$, on obtient $y=-x$. De $y+z=0$, on obtient $z=x$. Dans la troisième équation :
\[
x+z=2x=0,
\]
donc $x=0$, puis $y=0$ et $z=0$. Ainsi
\[
\Ker(u)=\{0\}.
\]
L'application est injective. Comme c'est un endomorphisme de $\R^3$, elle est bijective. Donc
\[
\rg(u)=3,\qquad \Img(u)=\R^3.
\]

\corr{16}
On a
\[
u(a,b)=a+aX+bX+bX^2.
\]
Dans la base $(1,X,X^2)$ :
\[
u(1,0)=1+X,\qquad u(0,1)=X+X^2.
\]
Donc
\[
\Mat(u)=
\begin{pmatrix}
1&0\\
1&1\\
0&1
\end{pmatrix}.
\]
Si $u(a,b)=0$, alors les coefficients de $1,X,X^2$ donnent
\[
a=0,\qquad a+b=0,\qquad b=0.
\]
Donc
\[
\Ker(u)=\{0\}.
\]
L'image est
\[
\Img(u)=\Vect(1+X,\;X+X^2).
\]
Elle est de dimension $2$. L'application est injective mais non surjective, car $\R_2[X]$ est de dimension $3$.

\corr{17}
$\varphi$ est une forme linéaire non nulle. Son noyau est défini par
\[
x+2y-z+t=0.
\]
On pose $y=a$, $z=b$, $t=c$. Alors
\[
x=-2a+b-c.
\]
Donc
\[
(x,y,z,t)=a(-2,1,0,0)+b(1,0,1,0)+c(-1,0,0,1).
\]
Ainsi
\[
\Ker(\varphi)=\Vect((-2,1,0,0),(1,0,1,0),(-1,0,0,1)).
\]
Cette famille est libre, donc elle est une base du noyau. Le noyau est de dimension $3$, c'est donc un hyperplan de $\R^4$.

\corr{18}
Soit $y\in\Img(u)$. Il existe $x\in E$ tel que $y=u(x)$. Comme $(e_1,\dots,e_n)$ engendre $E$, on peut écrire
\[
x=\lambda_1e_1+\cdots+\lambda_ne_n.
\]
Alors
\[
y=u(x)=\lambda_1u(e_1)+\cdots+\lambda_nu(e_n).
\]
Donc $(u(e_1),\dots,u(e_n))$ engendre $\Img(u)$. Un exemple où cette famille n'est pas libre : prendre $u:\R^2\to\R$, $u(x,y)=x$. L'image de la base canonique est $(1,0)$ dans $\R$, famille liée.

\corr{19}
Supposons
\[
\lambda_1u(e_1)+\cdots+\lambda_nu(e_n)=0.
\]
Par linéarité :
\[
u(\lambda_1e_1+\cdots+\lambda_ne_n)=0.
\]
Donc
\[
\lambda_1e_1+\cdots+\lambda_ne_n\in\Ker(u).
\]
Comme $u$ est injective, $\Ker(u)=\{0\}$. Ainsi
\[
\lambda_1e_1+\cdots+\lambda_ne_n=0.
\]
La famille $(e_1,\dots,e_n)$ étant libre, tous les $\lambda_i$ sont nuls. Donc la famille image est libre.

\corr{20}
Les matrices sont
\[
A=
\begin{pmatrix}
1&1\\
1&-1\\
0&2
\end{pmatrix},
\qquad
B=
\begin{pmatrix}
1&0&1\\
0&1&-1
\end{pmatrix}.
\]
Directement :
\[
u(x,y)=(x+y,x-y,2y),
\]
donc
\[
v(u(x,y))=(x+y+2y,\;x-y-2y)=(x+3y,\;x-3y).
\]
Par produit matriciel :
\[
BA=
\begin{pmatrix}
1&0&1\\
0&1&-1
\end{pmatrix}
\begin{pmatrix}
1&1\\
1&-1\\
0&2
\end{pmatrix}
=
\begin{pmatrix}
1&3\\
1&-3
\end{pmatrix}.
\]
C'est bien la matrice de $(x,y)\mapsto(x+3y,x-3y)$.

\subsection{Corrigé des exercices de démonstration}

\corr{21}
On écrit
\[
0_E=0_E+0_E.
\]
Donc
\[
u(0_E)=u(0_E+0_E)=u(0_E)+u(0_E).
\]
En soustrayant $u(0_E)$ aux deux membres, on obtient
\[
u(0_E)=0_F.
\]

\corr{22}
On a $u(0_E)=0_F$, donc $0_E\in\Ker(u)$.
Si $x,y\in\Ker(u)$, alors
\[
u(x+y)=u(x)+u(y)=0_F+0_F=0_F.
\]
Donc $x+y\in\Ker(u)$. Si $\lambda\in\R$ et $x\in\Ker(u)$, alors
\[
u(\lambda x)=\lambda u(x)=0_F.
\]
Donc $\lambda x\in\Ker(u)$. Le noyau est un sous-espace vectoriel.

\corr{23}
Comme $u(0_E)=0_F$, on a $0_F\in\Img(u)$. Soient $y_1,y_2\in\Img(u)$. Il existe $x_1,x_2\in E$ tels que
\[
y_1=u(x_1),\qquad y_2=u(x_2).
\]
Alors
\[
\lambda y_1+\mu y_2=\lambda u(x_1)+\mu u(x_2)=u(\lambda x_1+\mu x_2).
\]
Donc $\lambda y_1+\mu y_2\in\Img(u)$. Ainsi $\Img(u)$ est un sous-espace vectoriel.

\corr{24}
Si $u$ est injective et si $x\in\Ker(u)$, alors
\[
u(x)=0_F=u(0_E).
\]
Par injectivité, $x=0_E$. Donc $\Ker(u)=\{0_E\}$.

Réciproquement, supposons $\Ker(u)=\{0_E\}$. Si $u(x)=u(y)$, alors
\[
u(x-y)=u(x)-u(y)=0_F.
\]
Donc $x-y\in\Ker(u)$, puis $x-y=0_E$, et donc $x=y$. L'application est injective.

\corr{25}
Tout $x\in E$ s'écrit de manière unique
\[
x=\lambda_1e_1+\cdots+\lambda_ne_n.
\]
On définit
\[
u(x)=\lambda_1f_1+\cdots+\lambda_nf_n.
\]
Cette définition est correcte grâce à l'unicité des coordonnées dans une base. Elle est linéaire car les coordonnées d'une combinaison linéaire se combinent linéairement.
Si $v$ est une autre application linéaire telle que $v(e_i)=f_i$, alors pour
\[
x=\sum_{i=1}^n\lambda_i e_i,
\]
on a
\[
v(x)=\sum_{i=1}^n\lambda_i v(e_i)=\sum_{i=1}^n\lambda_i f_i=u(x).
\]
Donc $v=u$. L'application est unique.

\corr{26}
Soit $(e_1,\dots,e_k)$ une base de $\Ker(u)$, complétée en une base
\[
(e_1,\dots,e_k,e_{k+1},\dots,e_n)
\]
de $E$.
Montrons que $(u(e_{k+1}),\dots,u(e_n))$ engendre $\Img(u)$. Soit $y\in\Img(u)$. Il existe $x\in E$ tel que $y=u(x)$. Écrivons
\[
x=\sum_{i=1}^n \lambda_i e_i.
\]
Alors
\[
u(x)=\sum_{i=1}^n\lambda_i u(e_i)=\sum_{i=k+1}^n\lambda_i u(e_i),
\]
car $u(e_i)=0$ pour $1\leq i\leq k$.
Montrons maintenant la liberté. Si
\[
\sum_{i=k+1}^n\lambda_i u(e_i)=0,
\]
alors
\[
u\left(\sum_{i=k+1}^n\lambda_i e_i\right)=0.
\]
Donc
\[
\sum_{i=k+1}^n\lambda_i e_i\in\Ker(u)=\Vect(e_1,\dots,e_k).
\]
On obtient alors une relation linéaire entre les vecteurs d'une base de $E$, donc tous les coefficients $\lambda_i$ sont nuls. La famille est libre. Elle est donc une base de $\Img(u)$, avec $n-k$ vecteurs. Ainsi
\[
\rg(u)=n-k.
\]
Donc
\[
\dim(E)=n=k+(n-k)=\dim(\Ker(u))+\rg(u).
\]

\corr{27}
Une application linéaire $u:E\to F$ est déterminée par les images d'une base $(e_1,\dots,e_n)$ de $E$. Pour chaque $u(e_j)$, il faut $p$ coordonnées dans une base de $F$. Comme il y a $n$ vecteurs de base dans $E$, il faut donc $pn$ coefficients. Ainsi
\[
\dim(\mathcal{L}(E,F))=pn.
\]

\corr{28}
Pour $x,y\in E$ et $\lambda,\mu\in\R$ :
\[
(v\circ u)(\lambda x+\mu y)=v(u(\lambda x+\mu y)).
\]
Comme $u$ est linéaire :
\[
u(\lambda x+\mu y)=\lambda u(x)+\mu u(y).
\]
Puis, comme $v$ est linéaire :
\[
v(\lambda u(x)+\mu u(y))=\lambda v(u(x))+\mu v(u(y)).
\]
Donc
\[
(v\circ u)(\lambda x+\mu y)=\lambda(v\circ u)(x)+\mu(v\circ u)(y).
\]
Ainsi $v\circ u$ est linéaire.

\corr{29}
Soient $y_1,y_2\in F$. Posons
\[
x_1=u^{-1}(y_1),\qquad x_2=u^{-1}(y_2).
\]
Alors $u(x_1)=y_1$ et $u(x_2)=y_2$. Par linéarité :
\[
u(\lambda x_1+\mu x_2)=\lambda y_1+\mu y_2.
\]
En appliquant $u^{-1}$ :
\[
u^{-1}(\lambda y_1+\mu y_2)=\lambda x_1+\mu x_2
=\lambda u^{-1}(y_1)+\mu u^{-1}(y_2).
\]
Donc $u^{-1}$ est linéaire.

\corr{30}
Comme $\varphi$ est non nulle, il existe $x_0\in E$ tel que $\varphi(x_0)\neq 0$. L'image de $\varphi$ est un sous-espace non nul de $\R$, donc
\[
\Img(\varphi)=\R.
\]
Ainsi $\rg(\varphi)=1$. Par le théorème du rang :
\[
\dim(E)=\dim(\Ker(\varphi))+1.
\]
Donc
\[
\dim(\Ker(\varphi))=n-1.
\]
Ainsi $\Ker(\varphi)$ est un hyperplan.

\subsection{Corrigé des exercices de réflexion}

\corr{31}
Oui. Par exemple
\[
u:\R^2\to\R^3,\qquad u(x,y)=(x,y,0).
\]
Si $u(x,y)=(0,0,0)$, alors $x=0$ et $y=0$, donc $u$ est injective. Elle n'est pas surjective car $(0,0,1)$ n'a pas d'antécédent. Par dimension, une application de $\R^2$ vers $\R^3$ ne peut pas être surjective, car son rang est au plus $2$.

\corr{32}
Oui. Par exemple
\[
u:\R^3\to\R^2,\qquad u(x,y,z)=(x,y).
\]
Elle est surjective car tout $(a,b)\in\R^2$ est l'image de $(a,b,0)$. Elle n'est pas injective car
\[
\Ker(u)=\{(0,0,z)\mid z\in\R\}.
\]
Par le théorème du rang, son noyau est de dimension $1$.

\corr{33}
Soit $a=u(1)$. Pour tout $x\in\R$, on a
\[
x=x\cdot 1.
\]
Par linéarité :
\[
u(x)=u(x\cdot 1)=x u(1)=ax.
\]
Ainsi toutes les applications linéaires de $\R$ dans $\R$ sont les homothéties
\[
x\mapsto ax.
\]

\corr{34}
Comme $\rg(u)=1$, $\Img(u)$ est une droite. Choisissons $e\neq 0$ dans $\Img(u)$. Alors
\[
\Img(u)=\Vect(e).
\]
Donc pour tout $x\in E$, il existe un unique scalaire $\varphi(x)$ tel que
\[
u(x)=\varphi(x)e.
\]
Pour montrer que $\varphi$ est linéaire, on écrit
\[
u(x+y)=u(x)+u(y)=\varphi(x)e+\varphi(y)e=(\varphi(x)+\varphi(y))e.
\]
Mais aussi
\[
u(x+y)=\varphi(x+y)e.
\]
Comme $e\neq 0$ :
\[
\varphi(x+y)=\varphi(x)+\varphi(y).
\]
Même raisonnement pour les scalaires. Donc $\varphi$ est linéaire. Enfin :
\[
u(x)=0\Longleftrightarrow \varphi(x)e=0\Longleftrightarrow \varphi(x)=0.
\]
Donc
\[
\Ker(u)=\Ker(\varphi).
\]

\corr{35}
Pour tout $x\in E$ :
\[
x=(x-p(x))+p(x).
\]
On a
\[
p(x-p(x))=p(x)-p^2(x)=p(x)-p(x)=0,
\]
donc $x-p(x)\in\Ker(p)$. De plus $p(x)\in\Img(p)$ par définition. Donc
\[
E=\Ker(p)+\Img(p).
\]
Si $y\in\Ker(p)\cap\Img(p)$, alors $y=p(x)$ et $p(y)=0$. Mais
\[
p(y)=p(p(x))=p^2(x)=p(x)=y.
\]
Donc $y=0$. La somme est directe :
\[
E=\Ker(p)\oplus\Img(p).
\]

\corr{36}
On a
\[
(s-\Id)(s+\Id)=s^2-\Id=0.
\]
Pour tout $x\in E$ :
\[
x=\frac{x+s(x)}{2}+\frac{x-s(x)}{2}.
\]
Posons
\[
x_+=\frac{x+s(x)}2.
\]
Alors
\[
s(x_+)=\frac{s(x)+s^2(x)}2=\frac{s(x)+x}{2}=x_+,
\]
donc $x_+\in\Ker(s-\Id)$. Posons
\[
x_-=\frac{x-s(x)}2.
\]
Alors
\[
s(x_-)=\frac{s(x)-s^2(x)}2=\frac{s(x)-x}{2}=-x_-,
\]
donc $x_-\in\Ker(s+\Id)$.
L'intersection est réduite à $\{0\}$ : si $v$ est à la fois fixé et envoyé sur son opposé, alors $v=-v$, donc $v=0$. Ainsi
\[
E=\Ker(s-\Id)\oplus\Ker(s+\Id).
\]

\corr{37}
Si $u^2=0$, alors pour tout $x\in E$ :
\[
u(u(x))=0.
\]
Donc $u(x)\in\Ker(u)$, ce qui donne
\[
\Img(u)\subset\Ker(u).
\]
Ainsi
\[
\rg(u)\leq \dim(\Ker(u)).
\]
Par le théorème du rang :
\[
\dim(E)=\dim(\Ker(u))+\rg(u).
\]
Comme $\dim(\Ker(u))\geq \rg(u)$ :
\[
\dim(E)\geq 2\rg(u).
\]
Donc
\[
2\rg(u)\leq \dim(E).
\]

\corr{38}
Si $u$ est injective, alors $\Ker(u)=\{0\}$. Par le théorème du rang :
\[
\dim(E)=\rg(u).
\]
Comme $\dim(E)=\dim(F)$, on obtient
\[
\rg(u)=\dim(F),
\]
donc $\Img(u)=F$, c'est-à-dire $u$ est surjective. Réciproquement, si $u$ est surjective, alors $\rg(u)=\dim(F)=\dim(E)$. Le théorème du rang impose alors
\[
\dim(\Ker(u))=0.
\]
Donc $u$ est injective. En dimension finie égale :
\[
u\text{ injective}\Longleftrightarrow u\text{ surjective}\Longleftrightarrow u\text{ bijective}.
\]

\corr{39}
Comme $p$ est un projecteur :
\[
E=\Ker(p)\oplus\Img(p).
\]
Prenons une base $(k_1,\dots,k_r)$ de $\Ker(p)$ et une base $(i_1,\dots,i_s)$ de $\Img(p)$. La réunion est une base de $E$.
Dans cette base, $p(k_j)=0$ pour les vecteurs du noyau, et $p(i_j)=i_j$ pour les vecteurs de l'image. Donc la matrice de $p$ est
\[
\begin{pmatrix}
0_r&0\\
0&I_s
\end{pmatrix}
\]
si l'on place d'abord la base du noyau, puis celle de l'image.

\corr{40}
On a
\[
\Img(v\circ u)\subset \Img(v),
\]
car tout vecteur de $\Img(v\circ u)$ est de la forme $v(u(x))$. Donc
\[
\rg(v\circ u)\leq \rg(v).
\]
De plus,
\[
\Img(v\circ u)=v(\Img(u)).
\]
L'image par $v$ d'un sous-espace de dimension $\rg(u)$ ne peut avoir une dimension supérieure à $\rg(u)$. Donc
\[
\rg(v\circ u)\leq \rg(u).
\]
Finalement :
\[
\rg(v\circ u)\leq \min(\rg(u),\rg(v)).
\]

\subsection{Corrigé des exercices d'oraux}

\corr{41}
La linéarité est immédiate car chaque coordonnée est linéaire en $x,y,z$. Le noyau est donné par
\[
\begin{cases}
x+y=0,\\
y+z=0,\\
x+z=0.
\end{cases}
\]
De $x+y=0$, $y=-x$. De $y+z=0$, $z=x$. Dans la troisième équation :
\[
x+z=2x=0,
\]
donc $x=0$, puis $y=z=0$. Ainsi
\[
\Ker(u)=\{0\}.
\]
Comme $u:\R^3\to\R^3$ et que son noyau est nul, $u$ est injective, donc bijective. Par conséquent
\[
\Img(u)=\R^3,\qquad \rg(u)=3.
\]
L'image est donc tout l'espace $\R^3$.

\corr{42}
Si $p^2=p$, alors pour tout $y\in\Img(p)$, il existe $x$ tel que $y=p(x)$. Donc
\[
p(y)=p(p(x))=p^2(x)=p(x)=y.
\]
Ainsi $p$ fixe les vecteurs de son image. Comme montré précédemment :
\[
E=\Ker(p)\oplus\Img(p).
\]
Réciproquement, si $E=F\oplus G$, tout $x\in E$ s'écrit de manière unique $x=f+g$ avec $f\in F$, $g\in G$. On définit
\[
p(x)=f.
\]
Alors $p$ est linéaire, $p^2=p$, $\Img(p)=F$ et $\Ker(p)=G$. Dans une base adaptée, sa matrice est une matrice bloc avec un bloc identité sur $F$ et un bloc nul sur $G$.

\corr{43}
Si $s^2=\Id$, alors $s$ est bijective et $s^{-1}=s$. On pose
\[
x_+=\frac{x+s(x)}2,\qquad x_-=\frac{x-s(x)}2.
\]
Alors $x=x_++x_-$. De plus :
\[
s(x_+)=x_+,\qquad s(x_-)=-x_-.
\]
Donc
\[
x_+\in\Ker(s-\Id),\qquad x_-\in\Ker(s+\Id).
\]
L'intersection est nulle, donc
\[
E=\Ker(s-\Id)\oplus\Ker(s+\Id).
\]
Géométriquement, $s$ est une symétrie par rapport au sous-espace des vecteurs fixes, parallèlement au sous-espace des vecteurs changés en leurs opposés.

\corr{44}
Puisque $\rg(u)=1$, $\Img(u)$ est une droite. Prenons $f\neq0$ dans $\Img(u)$. Alors
\[
\Img(u)=\Vect(f).
\]
Pour tout $x\in E$, il existe un unique scalaire $\varphi(x)$ tel que
\[
u(x)=\varphi(x)f.
\]
La linéarité de $u$ entraîne celle de $\varphi$ par identification devant $f\neq0$. Le noyau est
\[
\Ker(u)=\Ker(\varphi).
\]
Réciproquement, si $u(x)=\varphi(x)f$ avec $\varphi\neq0$ et $f\neq0$, alors l'image est contenue dans $\Vect(f)$ et n'est pas nulle, donc elle est exactement $\Vect(f)$. Le rang vaut donc $1$.

\corr{45}
On résout
\[
\begin{cases}
x+2y+z=0,\\
2x+4y+2z=0,\\
x+y=0.
\end{cases}
\]
La deuxième équation est deux fois la première. De $x+y=0$, on a $x=-y$. Dans la première :
\[
-y+2y+z=0,
\]
donc
\[
y+z=0,\qquad z=-y.
\]
Ainsi
\[
(x,y,z)=y(-1,1,-1),
\]
et
\[
\Ker(A)=\Vect((-1,1,-1)).
\]
Donc $\dim(\Ker(A))=1$. Par le théorème du rang :
\[
\rg(A)=3-1=2.
\]
L'image est engendrée par les colonnes :
\[
C_1=(1,2,1),\quad C_2=(2,4,1),\quad C_3=(1,2,0).
\]
On peut garder par exemple $C_1$ et $C_2$, qui ne sont pas colinéaires :
\[
\Img(A)=\Vect((1,2,1),(2,4,1)).
\]

\corr{46}
Si $u^2=0$, alors $\Img(u)\subset\Ker(u)$. Donc
\[
\rg(u)\leq \dim(\Ker(u)).
\]
Par le théorème du rang :
\[
\dim(E)=\dim(\Ker(u))+\rg(u).
\]
Si $\dim(E)=3$, alors
\[
2\rg(u)\leq3,
\]
donc
\[
\rg(u)\leq1.
\]
Un exemple non nul dans $\R^2$ :
\[
u(x,y)=(y,0).
\]
Alors
\[
u^2(x,y)=u(y,0)=(0,0).
\]

\corr{47}
On a
\[
v_1=(1,1),\qquad v_2=(1,-1).
\]
Ils ne sont pas colinéaires, donc forment une base de $\R^2$. Calculons :
\[
u(v_1)=u(1,1)=(4,4)=4v_1,
\]
et
\[
u(v_2)=u(1,-1)=(2,-2)=2v_2.
\]
Donc, dans la base $(v_1,v_2)$ :
\[
\Mat(u)=
\begin{pmatrix}
4&0\\
0&2
\end{pmatrix}.
\]
Dans la base canonique, la matrice est
\[
\begin{pmatrix}
3&1\\
1&3
\end{pmatrix}.
\]
La base $(v_1,v_2)$ rend donc la matrice plus simple.

\corr{48}
On définit
\[
\varphi(x,y,z,t)=x-y+2z+t.
\]
Alors
\[
H=\Ker(\varphi),
\]
avec $\varphi$ forme linéaire non nulle. Résolvons :
\[
x-y+2z+t=0.
\]
On pose $y=a$, $z=b$, $t=c$. Alors
\[
x=a-2b-c.
\]
Donc
\[
(x,y,z,t)=a(1,1,0,0)+b(-2,0,1,0)+c(-1,0,0,1).
\]
Ainsi
\[
H=\Vect((1,1,0,0),(-2,0,1,0),(-1,0,0,1)).
\]
Cette famille est libre, donc c'est une base de $H$. Ainsi $\dim(H)=3$, donc $H$ est un hyperplan de $\R^4$.

\corr{49}
Si $x\in\Ker(u)$, alors
\[
u(x)=0.
\]
Donc
\[
(v\circ u)(x)=v(0)=0,
\]
d'où
\[
\Ker(u)\subset\Ker(v\circ u).
\]
Si $v$ est injective, alors
\[
(v\circ u)(x)=0 \Rightarrow v(u(x))=0 \Rightarrow u(x)=0,
\]
donc égalité. De plus, pour tout $x\in E$ :
\[
(v\circ u)(x)=v(u(x))\in\Img(v),
\]
donc
\[
\Img(v\circ u)\subset\Img(v).
\]
Exemple d'inclusion stricte : prendre $u=0$. Alors $\Img(v\circ u)=\{0\}$, tandis que $\Img(v)$ peut être non nulle.

\corr{50}
On choisit un supplémentaire $G$ de $\Ker(u)$ dans $E$ :
\[
E=\Ker(u)\oplus G.
\]
La restriction $u_{|G}$ est injective : si $x\in G$ et $u(x)=0$, alors $x\in G\cap\Ker(u)$, donc $x=0$.
De plus, tout vecteur de $\Img(u)$ est de la forme $u(x)$, avec $x=k+g$ où $k\in\Ker(u)$ et $g\in G$. Alors
\[
u(x)=u(k)+u(g)=u(g).
\]
Donc
\[
u(G)=\Img(u).
\]
La restriction
\[
u_{|G}:G\to\Img(u)
\]
est donc injective et surjective, donc c'est un isomorphisme. Ainsi
\[
\dim(G)=\dim(\Img(u))=\rg(u).
\]
Comme
\[
\dim(E)=\dim(\Ker(u))+\dim(G),
\]
on retrouve
\[
\dim(E)=\dim(\Ker(u))+\rg(u).
\]

\end{document}
